\subsection{\texorpdfstring{$\minilang$}{MiniLang} Static Semantics and Judgments in Twelf}
File \code{typing.elf} from~\cite{twelf-tutorial} contains the encoding of
$\minilang$'s static semantics in Twelf.
Line 5, ``\code{of :\ exp -> typ -> type.}'',
defines the relation \code{of} representing
the judgment $e : \tau$.
It defines \code{of} to be a function kind or
\emph{type family}~\cite{GHCTypeFamily,TwelfTypeFamily}.
Type families are functions that return types instead of terms.
The ``\code{of}'' relation can be thought of as (1) a function
that returns types when applied to \code{exp} and \code{typ} terms
or (2) as a set of types \emph{indexed} by \code{exp} and \code{typ} terms.

Judgments are represented in Twelf as
\emph{dependent types}~\cite{TwelfTypeJudge,advancetypes_ch02,Twelf-DependentTypes}.
The \code{of} type family returns dependent types of kind \code{type}.
Hence, function \code{of} returns judgments.
% Introduce dependent types
Dependent types are types that include terms as one of its components.
Example dependent types come from the set of $n$-dimensional vectors
of real numbers denoted \code{Vec($n$)}.
For instance, suppose \code{3} is a term representing the number $3$,
$\nats$ is a type representing the set of natural numbers,
and \code{3} is of type $\nats$.
Also, assume \code{Vec} is a function that takes in a natural
number $n$ as input and returns a type representing
the set of $n$-dimensional vectors.
Then, \code{Vec(3)} is a dependent type representing
$3$-dimensional vectors of real numbers.
A term of type \code{Vec(3)} is the vector $[7, 3, 4]$.
Furthermore, function \code{Vec} denotes a type family
indexed by natural numbers:
\code{Vec(0)}, \code{Vec(1)}, \code{Vec(2)}, \code{Vec(3)}, \ldots

% Back to Twelf
The \code{of} type family is indexed by \code{exp} and \code{typ} terms.
The dependent type ``\mbox{\code{of $e$ $\tau$}}'' represents
the judgment $e : \tau$.
For example, dependent type ``\mbox{\code{of (enat z) num}}''
represents the proposition or judgment $z :\ \code{num}$.
A proof or derivation of a judgment/type is represented by a term of that type.
An example proof will be presented later in this section.

Inference rules in Twelf are modeled as functions that return terms
of dependent types.
Consider the function term \code{of/nat} defined on line 7
of file \code{typing.elf}.
Strings that begin with capital letters in Twelf are interpreted
to be parameters.
The \code{N} parameter in ``\mbox{\code{of (enat N) num}}'' is passed
to the \code{enat} function term.
Since \code{enat} only takes in values of type
\code{nat}, so does \code{of/nat}.
Twelf is able to infer that the \code{N} variable must be
bound to a term of type \code{nat}.

At first, it seems the \code{of/nat} function returns types.
For example, it seems ``\code{of/nat z}'' returns
the dependent type ``\mbox{\code{of (enat z) num}}''.
However, ``\code{of/nat z}'' actually returns
a \emph{term} of type ``\mbox{\code{of (enat z) num}}''.
The term returned by ``\code{of/nat z}'' represents a derivation
or proof of the judgment $z :\ \code{num}$;
this judgment is represented by type ``\mbox{\code{of (enat z) num}}'',
which is the type of term ``\code{of/nat z}''.
Hence, the \code{of/nat} term corresponds to rule \code{T.1},
which is the rule that allows us to derive that any natural
number has type \code{num}.
Similarly, the \code{of/str} function term takes in any string $s$
and returns a derivation/term of a judgment/type stating that
string $s$ has type \code{string}.

Function term \code{of/add} models rule \code{T.4}.
Premises of inference rules are represented by inputs that must be
terms of a dependent type.
For example, function \code{of/add} takes in two
(explicit~\cite{Twelf-Implicit}) arguments of dependent types.\footnote{
Function \code{of/add} actually takes in four arguments.
The first two arguments are the \code{exp} terms \code{E1} and
\code{E2}.
However, these arguments are \emph{implicit} arguments and
typically do not need to be mentioned because they can be inferred
from the last two \emph{explicit} arguments of dependent types.
For example, since the dependent type ``\mbox{\code{of E1 num}}'' includes
\code{E1} as one if its components, Twelf can extract \code{E1}
from this type.
For more information on implicit and explicit arguments
see~\cite{Twelf-Implicit}.}
Variables \code{E1} and \code{E2} are bound to terms of type \code{exp}.
Dependent type ``\mbox{\code{of E1 num}}'' represents the judgment
$e_1 : \code{num}$.
A term of type ``\mbox{\code{of E1 num}}'' represents a derivation
of $e_1 : \code{num}$.
Hence, \code{of/add} takes in a derivation of $e_1 : \code{num}$
and a derivation of $e_2 : \code{num}$ and returns
a derivation/term of type ``\mbox{\code{of (add E1 E2) num}}'';
this type represents judgment \mbox{\code{+($e_1$; $e_2$) $:$ num}}.

\subsubsection{Higher-Order Judgments in Twelf}
Inference rules involving typing contexts
such as rule \code{T.6} are \emph{hypothetical judgments}~\cite{TwelfHypothetical}
or judgments that make use of hypothetical assumptions.
Hypothetical judgments can be encoded in Twelf using
\emph{higher-order judgments}~\cite{TwelfHigherOrder}
or judgments that contain other judgments.
Higher-order judgments are encoded in Twelf as \emph{higher-order types}.
Higher-order types are functions types, where one of their input types
is also a function type.
Encoding higher-order judgments as higher-order types is similar to
the technique of higher-order abstract syntax discussed in
Section~\ref{sec:twelf-hoas}, where syntactic terms with binders
(e.g. the body of the \code{let} expression) are encoded as
terms with holes or equivalently
functions that take in terms as input and return terms.
Higher-order judgments are encoded with function types that have
input types representing the hypothetical assumptions.

Higher-order types representing hypothetical judgments involve
\emph{pi-types}, denoted $\Pi x:S.T$.
A pi-type is a special kind of function type.
Terms of pi-types are \emph{pi-abstractions}.
First, a \emph{lambda-abstraction}, denoted ``$\lambda x:S.e$'',
is a function term of a more conventional kind of function type,
 $\ftype{S}{T}$.
A function of this type takes in a term of type $S$ and
returns a term of type $T$.
However, a pi-abstraction of pi-type $\Pi x:S.T$ would map
a term $s$ of type $S$ to a term of type $[s / x]T$.
That is, the return type of a pi-type can vary according
to the argument supplied.
Hence, if $x$ is (syntactically) a part of $T$ in the pi-type
$\Pi x:S.T$, then a term returned by a function of that
pi-type is a term of a dependent type, since the return type
$T$ depends on the input term $x$.
If $x$ is not a part of $T$ in $\Pi x:S.T$,
then we abbreviate $\Pi x:S.T$ as $\ftype{S}{T}$,
which also signals that a lambda-abstraction can be of
this type, since the return type does not involve the argument.
Lastly, the pi-type $\Pi x:S.T$ is represented in Twelf syntax
as \code{\{x:S\} T}.

A common Twelf coding convention is used in
the return type of the \code{of/let} function term
given on line 20 of the \code{typing.elf} file:
``\mbox{\code{of (let E1 ([x] E2 x)) T2}}''.
This code snippet could have been replaced with
the shorter snippet:
``\mbox{\code{of (let E1 E2) T2}}''.
The \code{E2} in the shorter snippet
already denotes a term
of function type ``\mbox{\code{exp -> exp}}''
(or equivalently, an \code{exp} with free variable
occurrences).
However, we applied a Twelf coding convention when
we replaced \code{E2} with the equivalent function
term \code{([x] E2 x)}.
The wrapper function \code{([x] E2 x)} is used 
just to make it easier to see that \code{E2} is
a function term that takes in a single argument.

The \code{of/let} term represents rule \code{T.6}
and is of a higher-order type:
\begin{verbatim}
of/let :  ({x: exp} of x T1 -> of (E2 x) T2) ->
          of E1 T1 ->
          of (let E1 ([x] E2 x)) T2
\end{verbatim}
The second argument type ``\mbox{\code{of E1 T1}}''
corresponds to the premise ``$e_1 : \tau_1$''.
The type of \code{of/let}'s first argument is the pi-type:
``\mbox{\code{\{x: exp\} of x T1 -> of (E2 x) T2}}'';
let $P$ denote this pi-type.
A pi-abstraction of type $P$ takes in two arguments:
The first is an \code{exp} term, which will be bound to
\code{x}.
The second is a term of type ``\mbox{\code{of x T1}}''
representing a proof that \code{x} has type \code{T1}.
Given two such terms,
a pi-abstraction of type $P$ should return
a derivation of ``\mbox{\code{of (E2 x) T2}}'',
where \code{(E2 x)} is the \code{E2} term with its hole
filled in by the \code{exp} term that is bound to \code{x}.

The pi-type $P$ represents the second premise
``\mbox{$\Gamma , x : \tau_1 \vdash e_2 : \tau_2$}''
of rule ~\code{T.6}.
Passing a derivation \code{dx} of type ``\mbox{\code{of x T1}}''
to a pi-abstraction simulates extending the typing context
with the assumption ``\mbox{$x : \tau_1$}''.
The derivation \code{dx} of type ``\mbox{\code{of x T1}}''
can be used in the body of the pi-abstraction to
return a term/proof of ``\mbox{\code{of (E2 x) T2}}''.

